\documentclass[../main.tex]{subfiles}

\graphicspath{{\subfix{../images/}}}

\begin{document}

\clearpage
\thispagestyle{empty}
\vspace*{7cm}
\subsection{Esercitazione 8 - 11/12/2025}
\vspace*{1.5cm}

\vspace*{1.5cm}
\subsubsection{Esercizio 1 - Processo casuale stazionario}
\label{ss81}

Dato un processo casuale stazionario $X(t) = A + \cos(2\pi t+ \Phi)$, dove $A$ e $\Phi$ sono due variabili casuali statisticamente indipendenti, che variano da realizzazione a realizzazione, con distribuzione uniforme negli intervalli $1  \leq  A \leq 3$, $ e - \pi \eq \Phi \leq + \Phi$, rispettivamente.
Si richiede di:

\begin{enumerate}
	\item Determinare se il processo è ergodico
	\item Determinare la densità spettrale di potenza del processo
	\item Determinare la potenza del processo $Y(t)$ ottenuto filtrando il processo $X(t)$ con un filtro con funzione di trasferimento $H(f) = \frac{1}{1+j 2\pi f}$

\end{enumerate}

\subsubsection*{Soluzione esercizio \ref{ss81}}

\vspace*{1.5cm}
\subsubsection{Esercizio 2 - }
\label{ss82}

\subsubsection*{Soluzione esercizio \ref{ss82}}

\vspace*{1.5cm}
\subsubsection{Esercizio 3 - }
\label{ss83}

\subsubsection*{Soluzione esercizio \ref{ss83}}

\vspace*{1.5cm}
\subsubsection{Esercizio 4 - }
\label{ss84}

\subsubsection*{Soluzione esercizio \ref{ss84}}

\vspace*{1.5cm}
\subsubsection{Esercizio 5 - }
\label{ss85}

\subsubsection*{Soluzione esercizio \ref{ss85}}

\vspace*{1.5cm}
\subsubsection{Esercizio 6 - }
\label{ss85}

\subsubsection*{Soluzione esercizio \ref{ss86}}


\vspace*{1.5cm}
\subsubsection{Esercizio 7 - }
\label{ss87}

\subsubsection*{Soluzione esercizio \ref{ss87}}



\subsection{Esercitazione 8 - Extra}

\subsubsection{Esercizio 1 - }

\subsubsection{Esercizio 1 - }




\end{document}